% GENERATED FILE. Edit canonical lessons, then run npm run generate:books.
% canonical-source-hash: fnv1a64:fe97cb41d80ee4cf
% canonical-lessons: KA-C272-cellu, KA-C272-suri, KA-C272-beresu, KA-C272-kalaku, KA-C272-huri

\chapter{To spill, To pour, To mix, To stir, To roast}
\label{ch:ka-to-spill-to-pour-to-mix-to-stir-to-roast}
\hlchaptermodality{272}

\begin{chapteropening}
\textbf{By the end of this chapter:} \emph{I can say to spill, to pour, to mix, to stir and to roast.}

Complete the last lesson of chapter 272: I can say to spill, to pour, to mix, to stir and to roast.
\end{chapteropening}

\section[\texorpdfstring{\kn{ಚೆಲ್ಲು} (cellu)}{ಚೆಲ್ಲು (cellu)}]{\kn{ಚೆಲ್ಲು} (cellu) — to spill}

\label{lesson:KA-C272-cellu}



\begin{quote}
Before the new one: say the Kannada for an organisation, then the Kannada for a government department.
\end{quote}

\subsection*{You'll want to know: \kn{ಚೆಲ್ಲು}}
\textbf{\kn{ಚೆಲ್ಲು}} — \emph{cellu} — \textquotedblleft{}to spill\textquotedblright{}.

\begin{grammarlens}[title={What you've built}]
One more word for this chapter.
\end{grammarlens}

\subsection*{Your turn}
\begin{itemize}
\raggedright
  \item \emph{Say it:} \emph{cellu}
  \item \emph{Say it:} \emph{cellu}, once more
  \item \emph{Say it:} say \emph{cellu}
  \item \emph{Recall:} say the Kannada for an organisation, then the Kannada for a government department, then say \emph{cellu} again
\end{itemize}

\begin{tcolorbox}[breakable,colback=teal!4,colframe=teal!35!black,title={Before you move on}]
What does \kn{ಚೆಲ್ಲು} mean? (\textquotedblleft{}To spill\textquotedblright{}.) Say it once more.
\end{tcolorbox}
\section[\texorpdfstring{\kn{ಸುರಿ} (suri)}{ಸುರಿ (suri)}]{\kn{ಸುರಿ} (suri) — to pour}

\label{lesson:KA-C272-suri}



\begin{quote}
Before the new one: say the Kannada for a government department, then the Kannada for to spill.
\end{quote}

\subsection*{You'll want to know: \kn{ಸುರಿ}}
\textbf{\kn{ಸುರಿ}} — \emph{suri} — \textquotedblleft{}to pour\textquotedblright{}.

\begin{grammarlens}[title={What you've built}]
One more word for this chapter.
\end{grammarlens}

\subsection*{Your turn}
\begin{itemize}
\raggedright
  \item \emph{Say it:} \emph{suri}
  \item \emph{Say it:} \emph{suri}, once more
  \item \emph{Say it:} say \emph{suri}
  \item \emph{Recall:} say the Kannada for a government department, then the Kannada for to spill, then say \emph{suri} again
\end{itemize}

\begin{tcolorbox}[breakable,colback=teal!4,colframe=teal!35!black,title={Before you move on}]
What does \kn{ಸುರಿ} mean? (\textquotedblleft{}To pour\textquotedblright{}.) Say it once more.
\end{tcolorbox}
\section[\texorpdfstring{\kn{ಬೆರೆಸು} (beresu)}{ಬೆರೆಸು (beresu)}]{\kn{ಬೆರೆಸು} (beresu) — to mix}

\label{lesson:KA-C272-beresu}



\begin{quote}
Before the new one: say the Kannada for to spill, then the Kannada for to pour.
\end{quote}

\subsection*{You'll want to know: \kn{ಬೆರೆಸು}}
\textbf{\kn{ಬೆರೆಸು}} — \emph{beresu} — \textquotedblleft{}to mix\textquotedblright{}.

\begin{grammarlens}[title={What you've built}]
One more word for this chapter.
\end{grammarlens}

\subsection*{Your turn}
\begin{itemize}
\raggedright
  \item \emph{Say it:} \emph{beresu}
  \item \emph{Say it:} \emph{beresu}, once more
  \item \emph{Say it:} say \emph{beresu}
  \item \emph{Recall:} say the Kannada for to spill, then the Kannada for to pour, then say \emph{beresu} again
\end{itemize}

\begin{tcolorbox}[breakable,colback=teal!4,colframe=teal!35!black,title={Before you move on}]
What does \kn{ಬೆರೆಸು} mean? (\textquotedblleft{}To mix\textquotedblright{}.) Say it once more.
\end{tcolorbox}
\section[\texorpdfstring{\kn{ಕಲಕು} (kalaku)}{ಕಲಕು (kalaku)}]{\kn{ಕಲಕು} (kalaku) — to stir}

\label{lesson:KA-C272-kalaku}



\begin{quote}
Before the new one: say the Kannada for to pour, then the Kannada for to mix.
\end{quote}

\subsection*{You'll want to know: \kn{ಕಲಕು}}
\textbf{\kn{ಕಲಕು}} — \emph{kalaku} — \textquotedblleft{}to stir\textquotedblright{}.

\begin{grammarlens}[title={What you've built}]
One more word for this chapter.
\end{grammarlens}

\subsection*{Your turn}
\begin{itemize}
\raggedright
  \item \emph{Say it:} \emph{kalaku}
  \item \emph{Say it:} \emph{kalaku}, once more
  \item \emph{Say it:} say \emph{kalaku}
  \item \emph{Recall:} say the Kannada for to pour, then the Kannada for to mix, then say \emph{kalaku} again
\end{itemize}

\begin{tcolorbox}[breakable,colback=teal!4,colframe=teal!35!black,title={Before you move on}]
What does \kn{ಕಲಕು} mean? (\textquotedblleft{}To stir\textquotedblright{}.) Say it once more.
\end{tcolorbox}
\section[\texorpdfstring{\kn{ಹುರಿ} (huri)}{ಹುರಿ (huri)}]{\kn{ಹುರಿ} (huri) — to roast, to fry}

\label{lesson:KA-C272-huri}



\begin{quote}
Before the new one: say the Kannada for to mix, then the Kannada for to stir.
\end{quote}

\subsection*{You'll want to know: \kn{ಹುರಿ}}
\textbf{\kn{ಹುರಿ}} — \emph{huri} — \textquotedblleft{}to roast, to fry\textquotedblright{}.

\begin{grammarlens}[title={What you've built}]
One more word for this chapter.
\end{grammarlens}

\subsection*{Your turn}
\begin{itemize}
\raggedright
  \item \emph{Say it:} \emph{huri}
  \item \emph{Say it:} \emph{huri}, once more
  \item \emph{Say it:} say \emph{huri}
  \item \emph{Recall:} say the Kannada for to mix, then the Kannada for to stir, then say \emph{huri} again
\end{itemize}

\begin{tcolorbox}[breakable,colback=teal!4,colframe=teal!35!black,title={Before you move on}]
What does \kn{ಹುರಿ} mean? (\textquotedblleft{}To roast, to fry\textquotedblright{}.) Say it once more.
\end{tcolorbox}
